% appP_register.tex -- generated by writing/tools/blueprint_gen.py 1.0 (level 3) on 2026-10-14 from sources a9a6d40381; do not edit
\providecommand{\bpNum}[1]{\hyperref[#1]{\ref*{#1}}}%
\providecommand{\bpNumA}[1]{\hyperref[#1]{\textcolor{orange!65!black}{\ref*{#1}}}}%
\providecommand{\bpPage}[1]{\pageref*{#1}}%
\providecommand{\bpLab}[1]{{\ttfamily #1}}%
\providecommand{\bpBrk}{\allowbreak}%
\providecommand{\bpT}[2]{\facttagbox[#1]{\scriptsize #2}}%
\providecommand{\bpU}[1]{{\color{gray!70!black}#1}}%

\section{假设登记表}\label{sec:appP-register}

本表由各单元的 \texttt{assumption} 环境与假设登记表自动生成，按所属单元的书序排列；“被使用于”一列由各陈述的 \texttt{\string\uses} 数据计算（直接使用者的编号），与环境中「被使用于」字段的比较见检查 C-P10。
\begingroup\scriptsize\setlength{\tabcolsep}{2.5pt}\renewcommand{\arraystretch}{1.06}
\begin{longtable}{@{}>{\raggedright\arraybackslash}p{\dimexpr(\textwidth-12\tabcolsep-0mm)*152/1000\relax}>{\raggedright\arraybackslash}p{\dimexpr(\textwidth-12\tabcolsep-0mm)*86/1000\relax}>{\raggedright\arraybackslash}p{\dimexpr(\textwidth-12\tabcolsep-0mm)*155/1000\relax}>{\raggedright\arraybackslash}p{\dimexpr(\textwidth-12\tabcolsep-0mm)*49/1000\relax}>{\raggedright\arraybackslash}p{\dimexpr(\textwidth-12\tabcolsep-0mm)*135/1000\relax}>{\raggedright\arraybackslash}p{\dimexpr(\textwidth-12\tabcolsep-0mm)*301/1000\relax}>{\raggedright\arraybackslash}p{\dimexpr(\textwidth-12\tabcolsep-0mm)*119/1000\relax}@{}}
\caption{编号假设登记表（由 \texttt{assumption} 环境生成）}\label{tab:appP-register}\\
\toprule
\textbf{编号／标签} & \textbf{单元} & \textbf{内容} & \textbf{标签} & \textbf{地位} & \textbf{被使用于} & \textbf{解除途径} \\
\midrule
\endfirsthead
\caption[]{编号假设登记表（由 \texttt{assumption} 环境生成）（续）}\\
\toprule
\textbf{编号／标签} & \textbf{单元} & \textbf{内容} & \textbf{标签} & \textbf{地位} & \textbf{被使用于} & \textbf{解除途径} \\
\midrule
\endhead
\midrule
\multicolumn{7}{r@{}}{\footnotesize（续下页）}\\
\endfoot
\bottomrule
\endlastfoot
\bpLab{\detokenize{asm:}\bpBrk \detokenize{intro-}\bpBrk \detokenize{wightman}} & \mbox{第~1~章} & QCD 的 Wightman 结构 (W1)--(W4) & \bpT{Plum}{C} & 全局非微扰假设。无全阶论证 & \hyperref[prop:prelim-optical]{3.24}；\bpBrk \hyperref[prop:ee-R-ImPi]{6.2}, \bpBrk \hyperref[prop:ee-analytic]{4}, \bpBrk \hyperref[prop:ee-adler-dispersion]{6}, \bpBrk \hyperref[prop:ee-contour]{17}, \bpBrk \hyperref[prop:ee-smearing]{24}；\bpBrk \hyperref[prop:pdf-invariance]{8.4}, \bpBrk \hyperref[prop:pdf-support-hadron]{7}, \bpBrk \hyperref[thm:pdf-moments-local]{13}, \bpBrk \hyperref[thm:pdf-dis-xspace]{22}, \bpBrk \hyperref[prop:pdf-sum-rules]{23}, \bpBrk \hyperref[prop:pdf-timelike-dglap]{29}；\bpBrk \hyperref[thm:heavy-inclusive-b]{20.6}；\bpBrk \hyperref[conj:power-renormalon-ht]{21.15}；\bpBrk \hyperref[thm:lamet-exp-decay]{22.32} & （单元尚未写入） \\
\bpLab{\detokenize{asm:}\bpBrk \detokenize{intro-}\bpBrk \detokenize{composite}} & \mbox{第~1~章} & 规范不变复合算符与光线算符的存在与重整化群 & \bpT{Plum}{C} & 全局非微扰假设。无全阶论证 & \hyperref[prop:pdf-invariance]{8.4}, \bpBrk \hyperref[prop:pdf-support-hadron]{7}, \bpBrk \hyperref[thm:pdf-moments-local]{13}, \bpBrk \hyperref[thm:pdf-dglap]{14}, \bpBrk \hyperref[thm:pdf-dis-xspace]{22}, \bpBrk \hyperref[prop:pdf-sum-rules]{23}, \bpBrk \hyperref[prop:pdf-timelike-dglap]{29}；\bpBrk \hyperref[prop:dy-scale]{12.22}；\bpBrk \hyperref[thm:heavy-inclusive-b]{20.6}；\bpBrk \hyperref[conj:power-renormalon-ht]{21.15} & （单元尚未写入） \\
\bpLab{\detokenize{asm:}\bpBrk \detokenize{intro-}\bpBrk \detokenize{np-}\bpBrk \detokenize{ope}} & \mbox{第~1~章} & 非微扰算符乘积展开 (O1)(O2) & \bpT{Plum}{C} & 全局非微扰假设。无全阶论证 & \hyperref[thm:pdf-dis-xspace]{8.22}；\bpBrk \hyperref[conj:power-renormalon-ht]{21.15} & （单元尚未写入） \\
\bpLab{\detokenize{asm:}\bpBrk \detokenize{intro-}\bpBrk \detokenize{af-}\bpBrk \detokenize{coeff}} & \mbox{第~1~章} & 系数的渐近自由行为 & \bpT{Plum}{C} & 全局非微扰假设。无全阶论证 & \hyperref[thm:pdf-dis-xspace]{8.22}；\bpBrk \hyperref[thm:heavy-inclusive-b]{20.6}；\bpBrk \hyperref[conj:power-renormalon-ht]{21.15} & （单元尚未写入） \\
\bpLab{\detokenize{asm:}\bpBrk \detokenize{intro-}\bpBrk \detokenize{growth}} & \mbox{第~1~章} & 前向振幅的增长条件 & \bpT{Plum}{C} & 全局非微扰假设。无全阶论证 & \hyperref[thm:pdf-dis-xspace]{8.22} & （单元尚未写入） \\
\bpLab{\detokenize{asm:}\bpBrk \detokenize{intro-}\bpBrk \detokenize{bridge}} & \mbox{第~1~章} & 从部分子微扰论到强子的桥接假设 & \bpT{Plum}{C} & 全局非微扰假设。无全阶论证 & \hyperref[thm:method-template]{9.42}；\bpBrk \hyperref[thm:disall-factorization]{10.11}, \bpBrk \hyperref[thm:disall-heavy]{18}, \bpBrk \hyperref[prop:disall-nuclear]{20}；\bpBrk \hyperref[thm:dy-factorization]{12.17}, \bpBrk \hyperref[thm:dy-colour-singlet]{25}, \bpBrk \hyperref[thm:dy-polarized]{26}, \bpBrk \hyperref[thm:dy-higher-twist]{29}；\bpBrk \hyperref[thm:scet-thrust]{15.16}；\bpBrk \hyperref[conj:jets-hadronization]{16.6}, \bpBrk \hyperref[thm:jets-ngl-factorization]{10}；\bpBrk \hyperref[thm:heavy-bdpi]{20.2}--\hyperref[prop:heavy-qcdf-lp]{4}, \bpBrk \hyperref[thm:heavy-inclusive-b]{6}, \bpBrk \hyperref[thm:heavy-nrqcd-decay]{9}, \bpBrk \hyperref[prop:heavy-quarkonium-fragmentation]{13}；\bpBrk \hyperref[prop:power-dis-twist4-x]{21.1}, \bpBrk \hyperref[prop:power-twist3-sidis]{3}, \bpBrk \hyperref[conj:power-twist3-pp]{4} & （单元尚未写入） \\
\bpNumA{asm:lamet-aux-renorm}\newline \bpLab{\detokenize{asm:}\bpBrk \detokenize{lamet-}\bpBrk \detokenize{aux-}\bpBrk \detokenize{renorm}} & \mbox{第~22~章} & 含类空辅助静态色源的 QCD 的可重整性 (AX1)--(AX3) & \bpT{RoyalBlue}{A} & 文献有全阶论证，未作为精确陈述证明 & \hyperref[prop:lattice-renormalizability]{22.9}, \bpBrk \hyperref[asm:lamet-ope]{\textcolor{orange!65!black}{19}}, \bpBrk \hyperref[asm:lamet-gap]{\textcolor{orange!65!black}{31}}, \bpBrk \hyperref[conj:lamet-decay-qcd]{33}, \bpBrk \hyperref[asm:lamet-staple]{\textcolor{orange!65!black}{53}}；\bpBrk \hyperref[thm:yzliu-soft-hqet]{23.3} & 见该假设的「解除途径」 \\
\bpNumA{asm:lamet-ope}\newline \bpLab{\detokenize{asm:}\bpBrk \detokenize{lamet-}\bpBrk \detokenize{ope}} & \mbox{第~22~章} & 类空 Wilson 线算符的短距离 OPE（重整化微扰论） (OP1)(OP2) & \bpT{RoyalBlue}{A} & 文献有全阶论证，未作为精确陈述证明 & \hyperref[thm:lamet-coordinate]{22.20}, \bpBrk \hyperref[thm:lamet-parton-momentum]{21}, \bpBrk \hyperref[prop:lamet-tmc]{38}, \bpBrk \hyperref[prop:lamet-sdf]{48}, \bpBrk \hyperref[thm:lamet-quasi-gpd]{51} & 见该假设的「解除途径」 \\
\bpNumA{asm:lamet-gap}\newline \bpLab{\detokenize{asm:}\bpBrk \detokenize{lamet-}\bpBrk \detokenize{gap}} & \mbox{第~22~章} & 谱隙与边界值正则性 (GP1)--(GP3) & \bpT{Plum}{C} & 无全阶论证 & \hyperref[thm:lamet-exp-decay]{22.32}, \bpBrk \hyperref[conj:lamet-decay-qcd]{33} & 在构造性或公理化框架中对所考虑的道证明 (GP1)(GP2)；(GP3) 需要对禁闭的控制，超出现有方法。 \\
\bpNumA{asm:lamet-uniform}\newline \bpLab{\detokenize{asm:}\bpBrk \detokenize{lamet-}\bpBrk \detokenize{uniform}} & \mbox{第~22~章} & 扭度展开对 Ioffe 时间的一致性 & \bpT{Plum}{C} & 无全阶论证 & \hyperref[thm:lamet-hadronic]{22.46} & 见该假设的「解除途径」 \\
\bpNumA{asm:lamet-twist4}\newline \bpLab{\detokenize{asm:}\bpBrk \detokenize{lamet-}\bpBrk \detokenize{twist4}} & \mbox{第~22~章} & 领头幂次修正的结构（扭度 4） (TW1)(TW2) & \bpT{Plum}{C} & 无全阶论证 & \hyperref[asm:lamet-large-z]{\textcolor{orange!65!black}{22.45}}, \bpBrk \hyperref[thm:lamet-hadronic]{46} & 对强子矩阵元、能控制 $z^2$ 依赖的扭度 4 OPE（非微扰 OPE）。 \\
\bpNumA{asm:lamet-large-z}\newline \bpLab{\detokenize{asm:}\bpBrk \detokenize{lamet-}\bpBrk \detokenize{large-}\bpBrk \detokenize{z}} & \mbox{第~22~章} & 大距离区域的控制 (LZ1)(LZ2) & \bpT{Plum}{C} & 无全阶论证 & \hyperref[thm:lamet-hadronic]{22.46} & 强子矩阵元的定量指数衰减，常数对 $P^z$ 一致。 \\
\bpNumA{asm:lamet-staple}\newline \bpLab{\detokenize{asm:}\bpBrk \detokenize{lamet-}\bpBrk \detokenize{staple}} & \mbox{第~22~章} & 订书钉形类空 Wilson 线算符的重整化与 $L\to\infty$ 极限 (ST1)--(ST3) & \bpT{RoyalBlue}{A} & 文献有全阶论证，未作为精确陈述证明 & \hyperref[def:lamet-quasi-tmd]{22.54}, \bpBrk \hyperref[thm:lamet-quasi-tmd]{59}；\bpBrk \hyperref[prop:yzliu-qtmd-powercount]{23.23} & 见该假设的「解除途径」 \\
\bpNumA{asm:lamet-soft-rapidity}\newline \bpLab{\detokenize{asm:}\bpBrk \detokenize{lamet-}\bpBrk \detokenize{soft-}\bpBrk \detokenize{rapidity}} & \mbox{第~22~章} & 离光锥软函数的大快度结构 (SR) & \bpT{RoyalBlue}{A} & 文献有全阶论证，未作为精确陈述证明 & \hyperref[def:lamet-reduced-soft]{22.57}, \bpBrk \hyperref[thm:lamet-quasi-tmd]{59} & 证明快度切片论证中的红外控制步骤（第~23~章）。 \\
\bpNumA{asm:lamet-order-limits}\newline \bpLab{\detokenize{asm:}\bpBrk \detokenize{lamet-}\bpBrk \detokenize{order-}\bpBrk \detokenize{limits}} & \mbox{第~22~章} & 极限次序 (OL1)(OL2) & \bpT{RoyalBlue}{A} & 文献有全阶论证，未作为精确陈述证明 & \hyperref[thm:lamet-quasi-tmd]{22.59} & 对 (OL1) 给出超出单圈的图形证明，并证明 (OL2)。 \\
\bpNumA{asm:yzliu-heavy-pair}\newline \bpLab{\detokenize{asm:}\bpBrk \detokenize{yzliu-}\bpBrk \detokenize{heavy-}\bpBrk \detokenize{pair}} & \mbox{第~23~章} & 静态色单态 $Q\bar Q$ 对的谱假设 (HP1)--(HP3) & \bpT{Plum}{C} & [C]——非微扰谱假设，(HP1) 有格点静态势作为证据。 & \hyperref[cor:yzliu-soft-form-factor]{23.6} & (HP1) 需要静态色单态扇区的能隙（禁闭），(HP2)(HP3) 需要非微扰 HQET 中静态源的 LSZ 型构造。 \\
\bpNumA{asm:yzliu-wightman-pt}\newline \bpLab{\detokenize{asm:}\bpBrk \detokenize{yzliu-}\bpBrk \detokenize{wightman-}\bpBrk \detokenize{pt}} & \mbox{第~23~章} & 维数正规化微扰论中的 Wightman 结构 (WP1)(WP2) & \bpT{RoyalBlue}{A} & 文献有全阶论证，未作为精确陈述证明 & \hyperref[thm:yzliu-analyticity]{23.9}, \bpBrk \hyperref[thm:yzliu-time-ordering]{20}, \bpBrk \hyperref[conj:yzliu-complex-wilson]{22}, \bpBrk \hyperref[thm:yzliu-jet-continuation]{31} & 见该假设的「解除途径」 \\
\bpNumA{asm:yzliu-lightcone-uniform}\newline \bpLab{\detokenize{asm:}\bpBrk \detokenize{yzliu-}\bpBrk \detokenize{lightcone-}\bpBrk \detokenize{uniform}} & \mbox{第~23~章} & 光锥极限对 ${\rm Im}\,\theta$ 的一致性 (UL) & \bpT{BurntOrange}{P} & 仅在特例、有限阶或模型中证明 & \hyperref[prop:yzliu-universality-classes]{23.16} & 在每一阶给出 $\order{\ee^{-\theta}}$ 余项在带形内的一致界。 \\
\bpLab{\detokenize{asm:}\bpBrk \detokenize{front-}\bpBrk \detokenize{mor-}\bpBrk \detokenize{complete}} & \mbox{第~24~章} & 区域集合的完备性 (MR) & \bpT{Plum}{C} & 无全阶论证 & — & （单元尚未写入） \\
\bpLab{\detokenize{asm:}\bpBrk \detokenize{front-}\bpBrk \detokenize{scet-}\bpBrk \detokenize{modes}} & \mbox{第~24~章} & SCET 的模式完备性与 Glauber 模式的缺席 (SM1)(SM2) & \bpT{RoyalBlue}{A} & 文献有全阶论证，未作为精确陈述证明 & — & （单元尚未写入） \\
\bpLab{\detokenize{asm:}\bpBrk \detokenize{appM-}\bpBrk \detokenize{remainder}} & \mbox{附录~M} & 冻结 $q$ 双局域正规乘积的动量空间余项估计 & \bpT{RoyalBlue}{A} & 文献有全阶论证，未作为精确陈述证明 & — & （单元尚未写入） \\
\bpLab{\detokenize{asm:}\bpBrk \detokenize{appM-}\bpBrk \detokenize{forward}} & \mbox{附录~M} & 前向在壳无质量部分子矩阵元的 OPE 余项 & \bpT{RoyalBlue}{A} & 文献有全阶论证，未作为精确陈述证明 & — & （单元尚未写入） \\
\bpLab{\detokenize{asm:}\bpBrk \detokenize{appM-}\bpBrk \detokenize{brst-}\bpBrk \detokenize{ope}} & \mbox{附录~M} & OPE 的 BRST 结构：带一个复合旁观者的欧氏展开 (E$\varepsilon$) (B2)(B1) & \bpT{RoyalBlue}{A} & 文献有全阶论证，未作为精确陈述证明 & — & （单元尚未写入） \\
\bpLab{\detokenize{asm:}\bpBrk \detokenize{appM-}\bpBrk \detokenize{oswi}} & \mbox{附录~M} & 前向在壳胶子振幅的 Ward 恒等式 & \bpT{RoyalBlue}{A} & 文献有全阶论证，未作为精确陈述证明 & — & （单元尚未写入） \\
\bpLab{\detokenize{asm:}\bpBrk \detokenize{appN-}\bpBrk \detokenize{contour}} & \mbox{附录~N} & (GA1) 围道 & \bpT{RoyalBlue}{A} & 文献有全阶论证，未作为精确陈述证明 & \hyperref[cor:method-removing-eps]{9.32}, \bpBrk \hyperref[thm:method-subtraction]{33}；\bpBrk \hyperref[prop:disall-two-routes]{10.2}, \bpBrk \hyperref[thm:disall-renorm-compat]{6} & （单元尚未写入） \\
\bpLab{\detokenize{asm:}\bpBrk \detokenize{appN-}\bpBrk \detokenize{envelope}} & \mbox{附录~N} & (GA3) 幂次计数包络 & \bpT{RoyalBlue}{A} & 文献有全阶论证，未作为精确陈述证明 & \hyperref[lem:method-local-accuracy]{9.12}, \bpBrk \hyperref[thm:method-template]{42}；\bpBrk \hyperref[prop:disall-two-routes]{10.2}, \bpBrk \hyperref[thm:disall-factorization]{11}；\bpBrk \hyperref[rem:appE-actual-bound]{E.2} & （单元尚未写入） \\
\bpLab{\detokenize{asm:}\bpBrk \detokenize{appN-}\bpBrk \detokenize{separation}} & \mbox{附录~N} & (GA4) 分离性 & \bpT{RoyalBlue}{A} & 文献有全阶论证，未作为精确陈述证明 & \hyperref[lem:method-local-accuracy]{9.12}, \bpBrk \hyperref[thm:method-template]{42}；\bpBrk \hyperref[prop:disall-two-routes]{10.2}, \bpBrk \hyperref[thm:disall-factorization]{11} & （单元尚未写入） \\
\bpLab{\detokenize{asm:}\bpBrk \detokenize{appN-}\bpBrk \detokenize{lce}} & \mbox{附录~N} & 光锥序表示与协变值的等同（LCE） (M1)(M2) & \bpT{RoyalBlue}{A} & 文献有全阶论证，未作为精确陈述证明 & — & （单元尚未写入） \\
\bpLab{\detokenize{asm:}\bpBrk \detokenize{appN-}\bpBrk \detokenize{values}} & \mbox{附录~N} & 发散链项的取值规则 (V1)(V2) & 无 & 约定 & \hyperref[thm:method-template]{9.42} & （单元尚未写入） \\
\bpLab{\detokenize{asm:}\bpBrk \detokenize{appN-}\bpBrk \detokenize{collins}} & \mbox{附录~N} & 一般图的 Collins 条件 (CS1)(CS3) & \bpT{RoyalBlue}{A} & 文献有全阶论证，未作为精确陈述证明 & \hyperref[thm:method-subtraction]{9.33}, \bpBrk \hyperref[thm:method-template]{42}；\bpBrk \hyperref[prop:disall-two-routes]{10.2}, \bpBrk \hyperref[prop:disall-h5]{10}, \bpBrk \hyperref[thm:disall-factorization]{11}；\bpBrk \hyperref[thm:dy-factorization]{12.17} & （单元尚未写入） \\
\bpLab{\detokenize{asm:}\bpBrk \detokenize{appN-}\bpBrk \detokenize{ward-}\bpBrk \detokenize{product}} & \mbox{附录~N} & Ward 恒等式的乘积结构与减除的重组 (WI1)(WI2) & \bpT{RoyalBlue}{A} & 文献有全阶论证，未作为精确陈述证明 & \hyperref[thm:method-template]{9.42}；\bpBrk \hyperref[prop:disall-factorized-subtraction]{10.5}, \bpBrk \hyperref[thm:disall-factorization]{11} & （单元尚未写入） \\
\bpLab{\detokenize{asm:}\bpBrk \detokenize{appN-}\bpBrk \detokenize{brst-}\bpBrk \detokenize{nonlocal}} & \mbox{附录~N} & 多定域光线算符的 BRST 分类 & \bpT{Plum}{C} & 具名缺口。无全阶论证 & — & （单元尚未写入） \\
\bpLab{\detokenize{asm:}\bpBrk \detokenize{appN-}\bpBrk \detokenize{cutsum}} & \mbox{附录~N} & 割图求和的公共表示 (CU1)--(CU3) & \bpT{RoyalBlue}{A} & 文献有全阶论证，未作为精确陈述证明 & \hyperref[prop:method-fsi]{9.37}, \bpBrk \hyperref[thm:method-template]{42}；\bpBrk \hyperref[prop:disall-two-routes]{10.2}, \bpBrk \hyperref[thm:disall-factorization]{11}；\bpBrk \hyperref[thm:dy-factorization]{12.17} & （单元尚未写入） \\
\bpLab{\detokenize{asm:}\bpBrk \detokenize{appN-}\bpBrk \detokenize{uv}} & \mbox{附录~N} & 被积函数层次的紫外重整化 (U1)(U2) & \bpT{RoyalBlue}{A} & 文献有全阶论证，未作为精确陈述证明 & \hyperref[thm:method-subtraction]{9.33}, \bpBrk \hyperref[thm:method-template]{42}；\bpBrk \hyperref[prop:disall-factorized-subtraction]{10.5}, \bpBrk \hyperref[thm:disall-factorization]{11} & （单元尚未写入） \\
\bpLab{\detokenize{asm:}\bpBrk \detokenize{appN-}\bpBrk \detokenize{multren}} & \mbox{附录~N} & 共线因子的乘法重整化（维数正规化路线） (Z1)(Z3) & \bpT{RoyalBlue}{A} & 文献有全阶论证，未作为精确陈述证明 & \hyperref[prop:pdf-timelike-dglap]{8.29}；\bpBrk \hyperref[thm:method-template]{9.42}；\bpBrk \hyperref[thm:disall-renorm-compat]{10.6}, \bpBrk \hyperref[prop:disall-h5]{10}, \bpBrk \hyperref[thm:disall-factorization]{11} & （单元尚未写入） \\
\bpLab{\detokenize{asm:}\bpBrk \detokenize{appN-}\bpBrk \detokenize{ir-}\bpBrk \detokenize{regulator}} & \mbox{附录~N} & 保持 Ward 恒等式的部分子靶红外正规化 (Z2) & \bpT{Plum}{C} & 无全阶论证 & \hyperref[thm:disall-renorm-compat]{10.6} & （单元尚未写入） \\
\bpLab{\detokenize{asm:}\bpBrk \detokenize{appN-}\bpBrk \detokenize{tmd}} & \mbox{附录~N} & TMD 输入 (T1)--(T4) & \bpT{RoyalBlue}{A} & 文献有全阶论证，未作为精确陈述证明 & — & （单元尚未写入） \\
\bpLab{\detokenize{asm:}\bpBrk \detokenize{appN-}\bpBrk \detokenize{endpoint}} & \mbox{附录~N} & 端点条件 (E1)--(E3) & \bpT{RoyalBlue}{A} & 文献有全阶论证，未作为精确陈述证明 & — & （单元尚未写入） \\
\bpLab{\detokenize{asm:}\bpBrk \detokenize{appN-}\bpBrk \detokenize{analytic-}\bpBrk \detokenize{reg}} & \mbox{附录~N} & 两圈图样积分的辅助解析正规化 (RA1)(RA2) & \bpT{RoyalBlue}{A} & 文献有全阶论证，未作为精确陈述证明 & — & （单元尚未写入） \\
\bpLab{\detokenize{asm:}\bpBrk \detokenize{appO-}\bpBrk \detokenize{lcopt}} & \mbox{附录~O} & 光锥序微扰论的适用性 & \bpT{RoyalBlue}{A} & 文献有全阶论证，未作为精确陈述证明 & \hyperref[prop:method-fsi]{9.37}；\bpBrk \hyperref[lem:dy-fsi]{12.7}, \bpBrk \hyperref[thm:dy-factorization]{17} & （单元尚未写入） \\
\bpLab{\detokenize{asm:}\bpBrk \detokenize{appO-}\bpBrk \detokenize{adjacent-}\bpBrk \detokenize{swap}} & \mbox{附录~O} & 相邻交换对称性（ADJ） (ADJ) & \bpT{BurntOrange}{P} & 仅在特例、有限阶或模型中证明 & — & （单元尚未写入） \\
\bpLab{\detokenize{asm:}\bpBrk \detokenize{appO-}\bpBrk \detokenize{real-}\bpBrk \detokenize{jets}} & \mbox{附录~O} & 旁观者喷注函数的实性（REAL） (REAL) & \bpT{BurntOrange}{P} & 仅在特例、有限阶或模型中证明 & — & （单元尚未写入） \\
\bpLab{\detokenize{asm:}\bpBrk \detokenize{appO-}\bpBrk \detokenize{collinear-}\bpBrk \detokenize{ends}} & \mbox{附录~O} & 共线端点的匹配 & \bpT{RoyalBlue}{A} & 文献有全阶论证，未作为精确陈述证明 & — & （单元尚未写入） \\
\bpLab{\detokenize{asm:}\bpBrk \detokenize{appO-}\bpBrk \detokenize{envelope}} & \mbox{附录~O} & 纤维积分误差的可积包络 & \bpT{RoyalBlue}{A} & 文献有全阶论证，未作为精确陈述证明 & \hyperref[prop:method-glauber-deformation]{9.40}；\bpBrk \hyperref[prop:dy-deformation]{12.9}, \bpBrk \hyperref[thm:dy-factorization]{17} & （单元尚未写入） \\
\bpLab{\detokenize{asm:}\bpBrk \detokenize{appO-}\bpBrk \detokenize{mixing}} & \mbox{附录~O} & Glauber--软混合与重合连接的完整图集相消 & \bpT{RoyalBlue}{A} & 文献有全阶论证，未作为精确陈述证明 & \hyperref[prop:method-glauber-deformation]{9.40}；\bpBrk \hyperref[prop:dy-deformation]{12.9}, \bpBrk \hyperref[thm:dy-factorization]{17} & （单元尚未写入） \\
\bpLab{\detokenize{asm:}\bpBrk \detokenize{appO-}\bpBrk \detokenize{deformation}} & \mbox{附录~O} & 一般构型、相容性、软角区与 H-发射端点 (CS1)(CS2) & \bpT{RoyalBlue}{A} & 文献有全阶论证，未作为精确陈述证明 & \hyperref[prop:method-glauber-deformation]{9.40}；\bpBrk \hyperref[prop:dy-deformation]{12.9}, \bpBrk \hyperref[thm:dy-factorization]{17} & （单元尚未写入） \\
\bpLab{\detokenize{asm:}\bpBrk \detokenize{appO-}\bpBrk \detokenize{rewritten-}\bpBrk \detokenize{cs}} & \mbox{附录~O} & 改写形式与族求和形式的 Collins 条件 (CS1)(CS2) & \bpT{RoyalBlue}{A} & 文献有全阶论证，未作为精确陈述证明 & \hyperref[prop:method-glauber-deformation]{9.40}；\bpBrk \hyperref[prop:dy-deformation]{12.9}, \bpBrk \hyperref[thm:dy-factorization]{17} & （单元尚未写入） \\
\end{longtable}
\endgroup
